\title{2022 年高考数学 全国乙卷-文科}
\date{}
\maketitle

\section*{一、选择题}
本题共12小题，每小题5分，共60分.在每小题给出的四个选项中，只有一项是符合题目要求的.

\begin{question}[points=5]
集合 $M = \{2, 4, 6, 8, 10\}$, $N = \{ x \mid -1 < x < 6 \}$, 则 $M \cap N =$
\begin{choices}
  \item $\{2, 4\}$
  \item $\{2, 4, 6\}$
  \item $\{2, 4, 6, 8\}$
  \item $\{2, 4, 6, 8, 10\}$
\end{choices}
\begin{solution}
\textbf{答案：}A

\textbf{分析：}根据集合的交集运算即可解出.

\textbf{详解：}因为 $M = \{2, 4, 6, 8, 10\}$, $N = \{ x \mid -1 < x < 6 \}$, 所以 $M \cap N = \{2, 4\}$.
\end{solution}
\end{question}

\begin{question}[points=5]
设 $(1+2i)a + b = 2i$, 其中 $a, b$ 为实数，则
\begin{choices}
  \item $a=1, b=-1$
  \item $a=1, b=1$
  \item $a=-1, b=1$
  \item $a=-1, b=-1$
\end{choices}
\begin{solution}
\textbf{答案：}A

\textbf{分析：}根据复数代数形式的运算法则以及复数相等的概念即可解出.

\textbf{详解：}因为 $a,b\in\mathbb{R}$, $(a+b)+2ai = 2i$, 所以 $a+b=0, 2a=2$, 解得 $a=1, b=-1$.
\end{solution}
\end{question}

\begin{question}[points=5]
已知向量 $\vec{a}=(2,1)$, $\vec{b}=(-2,4)$, 则 $|\vec{a} - \vec{b}| =$
\begin{choices}
  \item 2
  \item 3
  \item 4
  \item 5
\end{choices}
\begin{solution}
\textbf{答案：}D

\textbf{分析：}先求得 $\vec{a} - \vec{b}$, 然后求得 $|\vec{a} - \vec{b}|$.

\textbf{详解：}因为 $\vec{a} - \vec{b} = (2,1) - (-2,4) = (4,-3)$, 所以 $|\vec{a} - \vec{b}| = \sqrt{4^2 + (-3)^2} = 5$.
\end{solution}
\end{question}

\begin{question}[points=5]
分别统计了甲、乙两位同学16周的各周课外体育运动时长（单位：h），得如下茎叶图：



则下列结论中错误的是
\begin{choices}
  \item 甲同学周课外体育运动时长的样本中位数为7.4
  \item 乙同学周课外体育运动时长的样本平均数大于8
  \item 甲同学周课外体育运动时长大于8的概率的估计值大于0.4
  \item 乙同学周课外体育运动时长大于8的概率的估计值大于0.6
\end{choices}
\begin{solution}
\textbf{答案：}C

\textbf{分析：}结合茎叶图、中位数、平均数、古典概型等知识确定正确答案.

\textbf{详解：}对于A，甲同学样本中位数为 $\frac{7.3+7.5}{2}=7.4$, 正确. 对于B，乙同学平均数为 $8.50625>8$, 正确. 对于C，甲同学时长大于8的概率估计值为 $\frac{6}{16}=0.375<0.4$, 错误. 对于D，乙同学时长大于8的概率估计值为 $\frac{13}{16}=0.8125>0.6$, 正确.
\end{solution}
\end{question}

\begin{question}[points=5]
若 $x, y$ 满足约束条件 $\begin{cases} x+y \ge 2 \\ x+2y \le 4 \\ y \ge 0 \end{cases}$, 则 $z = 2x - y$ 的最大值是
\begin{choices}
  \item $-2$
  \item 4
  \item 8
  \item 12
\end{choices}
\begin{solution}
\textbf{答案：}C

\textbf{分析：}作出可行域，数形结合即可得解.

\textbf{详解：}由题意作出可行域，如图阴影部分所示. 转化目标函数 $z=2x-y$ 为 $y=2x-z$, 上下平移直线，当直线过点 $(4,0)$ 时，截距最小，$z$ 最大，所以 $z_{\max}=2\times4-0=8$.
\end{solution}
\end{question}

\begin{question}[points=5]
设 $F$ 为抛物线 $C: y^2=4x$ 的焦点，点 $A$ 在 $C$ 上，点 $B(3,0)$, 若 $|AF| = |BF|$, 则 $|AB| =$
\begin{choices}
  \item 2
  \item $2\sqrt{2}$
  \item 3
  \item $3\sqrt{2}$
\end{choices}
\begin{solution}
\textbf{答案：}B

\textbf{分析：}根据抛物线上的点到焦点和准线的距离相等，从而求得点 $A$ 的横坐标，进而求得点 $A$ 坐标，即可得到答案.

\textbf{详解：}由题意得 $F(1,0)$, 则 $|AF|=|BF|=2$, 即点 $A$ 到准线 $x=-1$ 的距离为2，所以点 $A$ 的横坐标为 $-1+2=1$, 不妨设 $A$ 在 $x$ 轴上方，代入得 $A(1,2)$, 所以 $|AB|=\sqrt{(3-1)^2+(0-2)^2}=2\sqrt{2}$.
\end{solution}
\end{question}

\begin{question}[points=5]
执行下边的程序框图，输出的 $n =$
\begin{choices}
  \item 3
  \item 4
  \item 5
  \item 6
\end{choices}
\begin{solution}
\textbf{答案：}B

\textbf{分析：}根据框图循环计算即可.

\textbf{详解：}执行第一次循环：$b=1+2=3$, $a=3-1=2$, $n=2$, $\left|\frac{b^2}{a^2}-2\right|=\left|\frac{9}{4}-2\right|=\frac{1}{4}>0.01$; 第二次：$b=3+4=7$, $a=7-2=5$, $n=3$, $\left|\frac{49}{25}-2\right|=\frac{1}{25}>0.01$; 第三次：$b=7+10=17$, $a=17-5=12$, $n=4$, $\left|\frac{289}{144}-2\right|=\frac{1}{144}<0.01$, 输出 $n=4$.
\end{solution}
\end{question}

\begin{question}[points=5]
如图是下列四个函数中的某个函数在区间 $[-3,3]$ 的大致图像，则该函数是
\begin{choices}
  \item $y=\frac{-x^3+3x}{x^2+1}$
  \item $y=\frac{x^3-x}{x^2+1}$
  \item $y=\frac{2x\cos x}{x^2+1}$
  \item $y=\frac{2\sin x}{x^2+1}$
\end{choices}
\begin{solution}
\textbf{答案：}A

\textbf{分析：}由函数图像的特征结合函数的性质逐项排除即可得解.

\textbf{详解：}设 $f(x)=\frac{x^3-x}{x^2+1}$, 则 $f(1)=0$, 排除B. 设 $h(x)=\frac{2x\cos x}{x^2+1}$, 当 $x\in(0,\frac{\pi}{2})$ 时 $0<\cos x<1$, 则 $h(x)<\frac{2x}{x^2+1}\le 1$, 排除C. 设 $g(x)=\frac{2\sin x}{x^2+1}$, 则 $g(3)=\frac{2\sin3}{10}>0$, 排除D. 故选A.
\end{solution}
\end{question}

\begin{question}[points=5]
在正方体 $ABCD-A_1B_1C_1D_1$ 中，$E,F$ 分别为 $AB, BC$ 的中点，则
\begin{choices}
  \item 平面 $B_1EF \perp$ 平面 $BDD_1$
  \item 平面 $B_1EF \perp$ 平面 $A_1BD$
  \item 平面 $B_1EF \parallel$ 平面 $A_1AC$
  \item 平面 $B_1EF \parallel$ 平面 $A_1C_1D$
\end{choices}
\begin{solution}
\textbf{答案：}A

\textbf{分析：}证明 $EF \perp$ 平面 $BDD_1$, 即可判断A；如图建系，分别求出法向量判断BCD.

\textbf{详解：}在正方体中，$EF \parallel AC$, $AC \perp BD$, 且 $DD_1 \perp$ 底面 $ABCD$, 所以 $EF \perp BD$, $EF \perp DD_1$, 故 $EF \perp$ 平面 $BDD_1$, 又 $EF \subset$ 平面 $B_1EF$, 所以平面 $B_1EF \perp$ 平面 $BDD_1$, A正确. 建系可验证其余错误.
\end{solution}
\end{question}

\begin{question}[points=5]
已知等比数列 $\{a_n\}$ 的前3项和为168，$a_2 - a_5 = 42$, 则 $a_6 =$
\begin{choices}
  \item 14
  \item 12
  \item 6
  \item 3
\end{choices}
\begin{solution}
\textbf{答案：}D

\textbf{分析：}设等比数列的公比为 $q$, 由题意求出首项与公比，再根据通项公式即可得解.

\textbf{详解：}设公比为 $q$, 由 $a_2-a_5=42$ 得 $q\neq1$. 则 $\begin{cases} \frac{a_1(1-q^3)}{1-q}=168 \\ a_1q - a_1q^4 = 42 \end{cases}$, 解得 $a_1=96, q=\frac12$, 所以 $a_6 = a_1q^5 = 3$.
\end{solution}
\end{question}

\begin{question}[points=5]
函数 $f(x) = \cos x + (x+1)\sin x + 1$ 在区间 $[0, 2\pi]$ 的最小值、最大值分别为
\begin{choices}
  \item $-\frac{\pi}{2}, \frac{\pi}{2}$
  \item $-\frac{3\pi}{2}, \frac{\pi}{2}$
  \item $-\frac{\pi}{2}, \frac{\pi}{2}+2$
  \item $-\frac{3\pi}{2}, \frac{\pi}{2}+2$
\end{choices}
\begin{solution}
\textbf{答案：}D

\textbf{分析：}利用导数求得单调区间，从而判断最值.

\textbf{详解：}$f'(x) = -\sin x + \sin x + (x+1)\cos x = (x+1)\cos x$, 在 $(0,\frac{\pi}{2})$ 和 $(\frac{3\pi}{2},2\pi)$ 上 $f'(x)>0$, $f(x)$ 单增；在 $(\frac{\pi}{2},\frac{3\pi}{2})$ 上 $f'(x)<0$, $f(x)$ 单减. 又 $f(0)=f(2\pi)=2$, $f(\frac{\pi}{2})=\frac{\pi}{2}+2$, $f(\frac{3\pi}{2}) = -\frac{3\pi}{2}$, 所以最小值为 $-\frac{3\pi}{2}$, 最大值为 $\frac{\pi}{2}+2$.
\end{solution}
\end{question}

\begin{question}[points=5]
已知球 $O$ 的半径为1，四棱锥的顶点为 $O$, 底面的四个顶点均在球 $O$ 的球面上，则当该四棱锥的体积最大时，其高为
\begin{choices}
  \item $\frac{1}{3}$
  \item $\frac{1}{2}$
  \item $\frac{\sqrt{3}}{3}$
  \item $\frac{\sqrt{2}}{2}$
\end{choices}
\begin{solution}
\textbf{答案：}C

\textbf{分析：}证明当四棱锥为正四棱锥时体积最大，利用基本不等式或导数求最值.

\textbf{详解：}设底面所在小圆半径为 $r$, 高为 $h$, 则 $r^2+h^2=1$. 底面为正方形时面积最大 $S=2r^2$, 体积 $V=\frac13\cdot2r^2\cdot h = \frac23 r^2 h \le \frac23 \sqrt{\frac{r^4+r^4+2h^2}{3}} = \frac{4\sqrt{3}}{27}$, 当且仅当 $r^2=2h^2$ 即 $h=\frac{\sqrt{3}}{3}$ 时取等.
\end{solution}
\end{question}

\section*{二、填空题}
本题共4小题，每小题5分，共20分.

\begin{question}[points=5]
记 $S_n$ 为等差数列 $\{a_n\}$ 的前 $n$ 项和. 若 $2S_3 = 3S_2 + 6$, 则公差 $d =$ .
\fillin[{$2$}]
\begin{solution}
\textbf{答案：}$2$

\textbf{分析：}转化条件为 $2(a_1+2d)=2a_1+d+6$, 即可得解.

\textbf{详解：}由 $2S_3=3S_2+6$ 得 $2(a_1+a_2+a_3)=3(a_1+a_2)+6$, 化简得 $2a_3=a_1+a_2+6$, 即 $2(a_1+2d)=2a_1+d+6$, 解得 $d=2$.
\end{solution}
\end{question}

\begin{question}[points=5]
从甲、乙等5名同学中随机选3名参加社区服务工作，则甲、乙都入选的概率为.
\fillin[{$\frac{3}{10}$}]
\begin{solution}
\textbf{答案：}$\frac{3}{10}$

\textbf{分析：}根据古典概型计算即可.

\textbf{详解：}从5名同学中选3名，共有 $C_5^3=10$ 种选法，甲、乙都入选的选法有 $C_3^1=3$ 种，所以概率为 $\frac{3}{10}$.
\end{solution}
\end{question}

\begin{question}[points=5]
过四点 $(0,0), (4,0), (-1,1), (4,2)$ 中的三点的一个圆的方程为.
\fillin[{$(x-2)^2+(y-3)^2=13$ 或 $(x-2)^2+(y-1)^2=5$ 或 $(x-\frac{4}{3})^2+(y-\frac{7}{3})^2=\frac{65}{9}$ 或 $(x-\frac{8}{5})^2+(y-1)^2=\frac{169}{25}$}]
\begin{solution}
\textbf{答案：}$(x-2)^2+(y-3)^2=13$ 或 $(x-2)^2+(y-1)^2=5$ 或 $(x-\frac{4}{3})^2+(y-\frac{7}{3})^2=\frac{65}{9}$ 或 $(x-\frac{8}{5})^2+(y-1)^2=\frac{169}{25}$

\textbf{分析：}法一：设圆的一般方程，根据所选点的坐标解方程组；法二：利用圆的几何性质求圆心和半径.

\textbf{详解：}答案不唯一，以上四个方程任写一个即可.
\end{solution}
\end{question}

\begin{question}[points=5]
若 $f(x) = \ln\left| a + \frac{1}{1-x} \right| + b$ 是奇函数，则 $a=$ , $b=$ .
\fillin[{$a=-\frac12$, $b=\ln2$}]
\begin{solution}
\textbf{答案：}$a=-\frac12$, $b=\ln2$

\textbf{分析：}根据奇函数的定义即可求出.

\textbf{详解：}由奇函数定义域关于原点对称，得 $a=-\frac12$, 再由 $f(0)=0$ 得 $b=\ln2$. 经检验，$f(x)=\ln\left|\frac{1+x}{1-x}\right|$ 满足奇函数.
\end{solution}
\end{question}

\section*{三、解答题}
共70分.解答应写出文字说明、证明过程或演算步骤.第17~21题为必考题，每个试题考生都必须作答.第22、23题为选考题，考生根据要求作答.

\begin{question}[points=12]
记 $\triangle ABC$ 的内角 $A,B,C$ 的对边分别为 $a,b,c$, 已知 $\sin C \sin(A-B) = \sin B \sin(C-A)$.

(1) 若 $A=2B$, 求 $C$;

(2) 证明：$2a^2 = b^2 + c^2$.
\begin{solution}
\textbf{答案：}(1) $C = \frac{5\pi}{8}$; (2) 证明见解析.

\textbf{分析：}(1) 由已知和 $A=2B$ 得 $\sin C = \sin(C-A)$, 结合三角形内角和求 $C$.
(2) 展开两角差正弦，利用正弦定理和余弦定理化简.

\textbf{详解：}(1) 由 $A=2B$ 代入得 $\sin C\sin B = \sin B\sin(C-A)$, 因为 $\sin B\neq0$, 所以 $\sin C = \sin(C-A)$. 因为 $0<C<\pi$, $0<C-A<\pi$, 且 $C\neq C-A$, 所以 $C+(C-A)=\pi$, 又 $A=2B$, $A+B+C=\pi$, 解得 $C=\frac{5\pi}{8}$.
(2) 展开得 $\sin C(\sin A\cos B-\cos A\sin B) = \sin B(\sin C\cos A-\cos C\sin A)$, 由正弦定理得 $ac\cos B - bc\cos A = bc\cos A - ab\cos C$. 由余弦定理，$\frac12(a^2+c^2-b^2)-\frac12(b^2+c^2-a^2) = \frac12(b^2+c^2-a^2)-\frac12(a^2+b^2-c^2)$, 化简得 $2a^2=b^2+c^2$.
\end{solution}
\end{question}

\begin{question}[points=12]
如图，四面体 $ABCD$ 中，$AD\perp CD$, $AD=CD$, $\angle ADB = \angle BDC$, $E$ 为 $AC$ 的中点.

(1) 证明：平面 $BED \perp$ 平面 $ACD$;

(2) 设 $AB=BD=2$, $\angle ACB=60^\circ$, 点 $F$ 在 $BD$ 上，当 $\triangle AFC$ 的面积最小时，求三棱锥 $F-ABC$ 的体积.
\begin{solution}
\textbf{答案：}(1) 证明见解析; (2) $\frac{3}{4}$.

\textbf{分析：}(1) 通过证明 $AC\perp$ 平面 $BED$ 来证得面面垂直.
(2) 判断三角形 $AFC$ 面积最小时 $F$ 的位置，然后求体积.

\textbf{详解：}(1) 因为 $AD=CD$, $E$ 为 $AC$ 中点，所以 $AC\perp DE$. 由 $\triangle ADB \cong \triangle CDB$ 得 $AB=CB$, 所以 $AC\perp BD$. 又 $DE\cap BD=D$, 故 $AC\perp$ 平面 $BED$, 而 $AC\subset$ 平面 $ACD$, 所以平面 $BED \perp$ 平面 $ACD$.
(2) 由 (1) 及已知得 $\triangle ABC$ 是边长为2的等边三角形，$BE=\sqrt{3}$, $DE=1$, $DE\perp BE$. 由 $\triangle ADB\cong\triangle CDB$ 知 $AF=CF$, 所以 $EF\perp AC$. 当 $EF\perp BD$ 时 $\triangle AFC$ 面积最小，在 $\triangle BED$ 中，$EF=\frac{BE\cdot DE}{BD}=\frac{\sqrt{3}}{2}$, $BF=\frac{3}{2}$, 过 $F$ 作 $FH\perp BE$ 于 $H$, 则 $FH\parallel DE$, $FH=\frac{3}{4}$, 且 $FH\perp$ 平面 $ABC$, 所以 $V_{F-ABC}=\frac13\cdot S_{\triangle ABC}\cdot FH = \frac13\cdot\frac{\sqrt{3}}{4}\cdot4\cdot\frac34 = \frac{\sqrt{3}}{4}$.
\end{solution}
\end{question}

\begin{question}[points=12]
某地经过多年的环境治理，已将荒山改造成了绿水青山．为估计一林区某种树木的总材积量，随机选取了10棵这种树木，测量每棵树的根部横截面积（单位：$m^2$）和材积量（单位：$m^3$），得到如下数据：

\begin{tabular}{c|cccccccccc|c}

样本号 $i$ & 1 & 2 & 3 & 4 & 5 & 6 & 7 & 8 & 9 & 10 & 总和 \\ \hline

根部横截面积 $x_i$ & 0.04 & 0.06 & 0.04 & 0.08 & 0.08 & 0.05 & 0.05 & 0.07 & 0.07 & 0.06 & 0.6 \\

材积量 $y_i$ & 0.25 & 0.40 & 0.22 & 0.54 & 0.51 & 0.34 & 0.36 & 0.46 & 0.42 & 0.40 & 3.9 \\

\end{tabular}

并计算得 $\sum_{i=1}^{10} x_i^2 = 0.038$, $\sum_{i=1}^{10} y_i^2 = 1.6158$, $\sum_{i=1}^{10} x_i y_i = 0.2474$.

(1) 估计该林区这种树木平均一棵的根部横截面积与平均一棵的材积量；

(2) 求该林区这种树木的根部横截面积与材积量的样本相关系数（精确到0.01）；

(3) 现测量了该林区所有这种树木的根部横截面积，并得到所有这种树木的根部横截面积总和为 $186 m^2$．已知树木的材积量与其根部横截面积近似成正比．利用以上数据给出该林区这种树木的总材积量的估计值．

附：相关系数 $r = \frac{\sum_{i=1}^{n}(x_i-\bar{x})(y_i-\bar{y})}{\sqrt{\sum_{i=1}^{n}(x_i-\bar{x})^2 \sum_{i=1}^{n}(y_i-\bar{y})^2}}$, $\sqrt{1.896}\approx1.377$.
\begin{solution}
\textbf{答案：}(1) $0.06m^2$, $0.39m^3$; (2) $0.97$; (3) $1209m^3$.

\textbf{分析：}计算样本平均值得估计值；代入相关系数公式计算；利用正比例关系求总材积量.

\textbf{详解：}(1) $\bar{x}=0.06$, $\bar{y}=0.39$.
(2) $r = \frac{0.2474-10\times0.06\times0.39}{\sqrt{(0.038-10\times0.06^2)(1.6158-10\times0.39^2)}} = \frac{0.0134}{\sqrt{0.0001896}} \approx \frac{0.0134}{0.01377} \approx 0.97$.
(3) 设总材积量为 $Y$, 由 $\frac{0.06}{0.39}=\frac{186}{Y}$ 得 $Y=1209m^3$.
\end{solution}
\end{question}

\begin{question}[points=12]
已知函数 $f(x) = ax - \frac{1}{x} - (a+1)\ln x$.

(1) 当 $a=0$ 时，求 $f(x)$ 的最大值；

(2) 若 $f(x)$ 恰有一个零点，求 $a$ 的取值范围.
\begin{solution}
\textbf{答案：}(1) $-1$; (2) $(0, +\infty)$.

\textbf{分析：}(1) 求导得单调性，即可得最大值.
(2) 求导分类讨论 $a\le0$, $0<a<1$, $a=1$, $a>1$ 时函数的单调性和零点情况.

\textbf{详解：}(1) $a=0$ 时 $f(x)=-\frac1x-\ln x$, $f'(x)=\frac{1-x}{x^2}$, 当 $x\in(0,1)$ 时 $f'(x)>0$, $f(x)$ 单增；当 $x\in(1,+\infty)$ 时 $f'(x)<0$, $f(x)$ 单减，所以 $f(x)_{\max}=f(1)=-1$.
(2) $f'(x)=\frac{(ax-1)(x-1)}{x^2}$. 当 $a\le0$ 时，$ax-1<0$, $f(x)$ 在 $(0,1)$ 单增，在 $(1,+\infty)$ 单减，$f(1)=a-1<0$, 无零点. 当 $0<a<1$ 时，$f(x)$ 在 $(0,1), (\frac1a,+\infty)$ 单增，在 $(1,\frac1a)$ 单减，$f(1)=a-1<0$, 取 $m=(\frac3a+2)^2$, 则 $f(m)>0$, 所以仅有一个零点. 当 $a=1$ 时，$f'(x)\ge0$, $f(x)$ 单增，$f(1)=0$, 唯一零点. 当 $a>1$ 时，$f(x)$ 在 $(0,\frac1a), (1,+\infty)$ 单增，在 $(\frac1a,1)$ 单减，$f(1)=a-1>0$, 取 $n=\frac1{4(a+1)}$, 则 $f(n)<0$, 所以有一个零点. 综上 $a>0$.
\end{solution}
\end{question}

\begin{question}[points=12]
已知椭圆 $E$ 的中心为坐标原点，对称轴为 $x$ 轴、$y$ 轴，且过 $A(0,-2)$, $B(\frac32,-1)$ 两点.

(1) 求 $E$ 的方程；

(2) 设过点 $P(1,-2)$ 的直线交 $E$ 于 $M,N$ 两点，过 $M$ 且平行于 $x$ 轴的直线与线段 $AB$ 交于点 $T$, 点 $H$ 满足 $\overrightarrow{MT} = \overrightarrow{TH}$. 证明：直线 $HN$ 过定点.
\begin{solution}
\textbf{答案：}(1) $\frac{y^2}{4}+\frac{x^2}{3}=1$; (2) $(0,-2)$.

\textbf{分析：}(1) 设椭圆方程 $mx^2+ny^2=1$ 代入坐标求解.
(2) 设直线方程联立，表示出 $H$ 点坐标，证明 $HN$ 过定点.

\textbf{详解：}(1) 设 $mx^2+ny^2=1$, 代入得 $\begin{cases} 4n=1 \\ \frac94 m + n = 1 \end{cases}$, 解得 $m=\frac13$, $n=\frac14$, 所以 $\frac{y^2}{4}+\frac{x^2}{3}=1$.
(2) 直线 $AB: y=\frac23 x -2$. ①斜率不存在：$x=1$, 得 $M(1,-\frac{2\sqrt6}{3})$, $N(1,\frac{2\sqrt6}{3})$, $T(-\sqrt6+3,-\frac{2\sqrt6}{3})$, $H(-2\sqrt6+5,-\frac{2\sqrt6}{3})$, 直线 $HN: y=(2+\frac{2\sqrt6}{3})x-2$, 过 $(0,-2)$. ②斜率存在：设 $kx-y-(k+2)=0$, 联立得 $(3k^2+4)x^2-6k(2+k)x+3k(k+4)=0$. 由 $y=y_1$ 与 $AB$ 联立得 $T(\frac{3y_1}{2}+3, y_1)$, 由 $\overrightarrow{MT}=\overrightarrow{TH}$ 得 $H(3y_1+6-x_1, y_1)$. 写出 $HN$ 方程，代入 $(0,-2)$ 整理得恒等式，故过定点 $(0,-2)$.
\end{solution}
\end{question}

\section*{（二）选考题}
共10分.请考生在第22、23题中选定一题作答.多答按所答第一题评分.

\begin{question}[points=10]
[选修4-4：坐标系与参数方程]

在直角坐标系 $xOy$ 中，曲线 $C$ 的参数方程为 $\begin{cases} x = \sqrt{3}\cos 2t \\ y = 2\sin t \end{cases}$（$t$ 为参数），以坐标原点为极点，$x$ 轴正半轴为极轴建立极坐标系，已知直线 $l$ 的极坐标方程为 $\rho\sin(\theta+\frac{\pi}{3})+m=0$.

(1) 写出 $l$ 的直角坐标方程；

(2) 若 $l$ 与 $C$ 有公共点，求 $m$ 的取值范围.
\begin{solution}
\textbf{答案：}(1) $\sqrt{3}x+y+2m=0$; (2) $[-\frac{19}{12}, \frac{5}{2}]$.

\textbf{分析：}(1) 利用极坐标与直角坐标互化.
(2) 联立 $l$ 与 $C$ 的方程，利用参数方程转化为三角函数有解.

\textbf{详解：}(1) 由 $\rho\sin(\theta+\frac{\pi}{3})+m=0$ 得 $\frac12\rho\sin\theta+\frac{\sqrt3}{2}\rho\cos\theta+m=0$, 即 $\frac12 y+\frac{\sqrt3}{2}x+m=0$, 所以 $\sqrt3 x+y+2m=0$.
(2) 将 $x=\sqrt3\cos2t$, $y=2\sin t$ 代入得 $\sqrt3\cdot\sqrt3\cos2t + 2\sin t + 2m=0$, 即 $3(1-2\sin^2t)+2\sin t+2m=0$, 整理得 $2m=6\sin^2t -2\sin t -3$. 令 $u=\sin t\in[-1,1]$, 则 $2m=6u^2-2u-3$, $u\in[-1,1]$. 该二次函数在 $u=\frac16$ 取最小值 $-\frac{19}{6}$, 在 $u=-1$ 取最大值5, 所以 $2m\in[-\frac{19}{6},5]$, 即 $m\in[-\frac{19}{12},\frac52]$.
\end{solution}
\end{question}

\begin{question}[points=10]
[选修4-5：不等式选讲]

已知 $a,b,c$ 都是正数，且 $a^{\frac32}+b^{\frac32}+c^{\frac32}=1$, 证明：

(1) $abc \le \frac19$;

(2) $\frac{a}{b+c}+\frac{b}{a+c}+\frac{c}{a+b} \le \frac{1}{2\sqrt{abc}}$.
\begin{solution}
\textbf{答案：}证明见解析.

\textbf{分析：}(1) 利用三元均值不等式.
(2) 利用基本不等式放缩.

\textbf{详解：}(1) 由 $a^{\frac32}+b^{\frac32}+c^{\frac32} \ge 3\sqrt[3]{a^{\frac32}b^{\frac32}c^{\frac32}} = 3(abc)^{\frac12}$ 得 $1 \ge 3(abc)^{\frac12}$, 所以 $abc \le \frac19$.
(2) 由 $b+c\ge2\sqrt{bc}$ 得 $\frac{a}{b+c} \le \frac{a}{2\sqrt{bc}} = \frac{a^{\frac32}}{2\sqrt{abc}}$, 同理 $\frac{b}{a+c} \le \frac{b^{\frac32}}{2\sqrt{abc}}$, $\frac{c}{a+b} \le \frac{c^{\frac32}}{2\sqrt{abc}}$, 相加得 $\frac{a}{b+c}+\frac{b}{a+c}+\frac{c}{a+b} \le \frac{a^{\frac32}+b^{\frac32}+c^{\frac32}}{2\sqrt{abc}} = \frac{1}{2\sqrt{abc}}$.
\end{solution}
\end{question}

